\chapter{Conclusions and Future Work}
\label{ch:ch8_concl}

\section{Conclusions}

The work set out to design a key-generation engine for an integrated
circuit whose entropy comes from ring oscillators simulated at transistor
level, to turn their jitter into a bound through a stochastic model, to
integrate the source in an \ac{FPGA} with an ESP32, and to compare the
random sequences of both. The conclusions follow the objectives of
Section~\ref{sec:objectives}; each one states what was shown and on what
evidence, \meas{}, \simu{} or \est{}.

\paragraph{The entropy source can be dimensioned from the transistor
(O1--O3).} The five-stage ring on the open IHP SG13G2 process oscillates
at 2.73~GHz with a period jitter of $176.7 \pm 2.0$~fs \simu{} in the
typical corner at 27~\textdegree C, with the injected noise taken from
the device models themselves rather than assumed. The simulator cannot fix the
realization of its noise source, so the uncertainty comes from repeating
the simulation, and the bench resolution was stated before any figure was
read. The exact wrapped-Gaussian computation, not the truncated series
usually quoted, turns the jitter into a min-entropy per bit, and the
target $H_\infty \ge 0.98$ fixes $Q \ge 0.2286$. Across corners, supply
($\pm 10$\,\%) and mismatch, the worst case is the slow corner at
$-40$~\textdegree C and 1.08~V, where a five-stage ring would need a
fixed divider of $2^{23}$ \simu{}. The divider therefore follows from the
cold, slow, low-supply corner, not from the typical one, which alone
would have suggested $2^{20}$.

\paragraph{The sampler, not the ring, fixed the length of the ring.} The
flip-flop of the library that samples the ring captures rising and
falling inputs at different instants, and that biases the bit. A
criterion was written before the Monte Carlo that would test it: declare
no more entropy than the $\mu - 3\sigma$ point of the worst corner
supports, and if that falls below 0.60 correct the circuit instead of
lowering the figure. With the five-stage ring it fell to 0.532 \simu{}.
The sampled ring was lengthened to eleven stages and the Monte Carlo
repeated on 24 dies in each of five corners: 0.757 \simu{}, so the
declared entropy of 0.60 stands, with a 512-bit seed. The final source
has rings of eleven and thirteen stages and a minimum divider of \kd{25},
twice what the simulation of its worst corner asks for, and delivers 20
to 50~bit/s \est{}. That is a hundred times slower than the first
estimates, and it is the honest cost of the bias, the worst corner and
the margin together.

\paragraph{A battery of statistical estimators cannot replace the model.}
Bits drawn from the model with the IHP jitter, and ideal uniform bits,
were passed through the SP~800-90B tool along the same path as the ESP32
captures. At $10^6$ bits the tool saturates at $0.839 \pm 0.021$
\simu{}: every point of the model with $H_\infty \ge 0.97$ reads alike,
and with little entropy the tool overestimates. With longer sequences the
ceiling rises for every source, ESP32 included, and the tool still does
not separate them (Section~\ref{sec:sim_90b}). The board confirmed it
\meas{}: on the \ac{FPGA} the tool gives 0.82 to 0.88 for every divider
from $2^{12}$ to $2^{19}$, while the model with the measured jitter goes
from almost nothing to 0.9996 over that range. Statistical tests detect
defects; they do not certify entropy. For a ring oscillator the bound has
to come from the physics, and that is the argument for the stochastic
model.

\paragraph{The digital engine is correct against independent references
(O4).} The engine --- \ac{ERO} sampler, health tests with start-up test,
CTR\_DRBG with derivation function, AES-CMAC, zeroization and an I2C
register map --- is portable VHDL-2008. Its thirteen testbenches
reproduce the official FIPS-197, SP~800-38B and RFC~4493 vectors and the
outputs of a reference model validated on the \ac{CAVP} vectors, and the
security rules of the sequencer are checked as \ac{PSL} properties over
the whole simulation, not only in the steps that exercise them. Three
independent reviews found defects that the testbenches had not caught ---
output still enabled with a health alarm raised, the seed reaching the
capture buffer, intermediate values surviving zeroization, a divider
setting out of range accepted in silence ---; all were fixed and each now
has a test. The lesson is methodological: a design
verified by its own author against its own testbench is not verified.

\paragraph{The engine fits a small integrated circuit (O7).} Synthesized
on IHP SG13G2 standard cells, the engine occupies 95.5~kGE with one shared
32-bit AES core \simu{}, of which 48.7~kGE are the digital engine
without the capture buffer. The two optimizations applied --- one AES core
shared by the generator and the authenticator, then a 32-bit datapath ---
saved 23~kGE when they were made; the 13 to 25~kGE that this work first
estimated for an engine of this kind are not reached, and the path to
about 45~kGE is identified.

\paragraph{The ESP32 generator is not a reference to be trusted blindly
(O8, ESP32 side).} Thirty-eight captures from three chips under four radio
and \ac{ADC} conditions show no measurable effect of condition or chip
under SP~800-90B \meas{}. The restart test shows what the statistics
hide: without the noise that the bootloader injects, what the generator
produces after every reset is a window of one fixed sequence, and two
different chips produce the same one \meas{}. Its output is as good as its seeding, which
depends on the radio, the \ac{ADC} or the bootloader being active.

% PENDIENTE DE MARIO: KAT en placa (reformular O5 o anotar)
\paragraph{The whole engine works on an FPGA with an independent verifier
(O5).} The engine was implemented complete on a Basys~3 and driven by an
ESP32 over \ac{I2C}. In three sessions, on three different
implementations, the ESP32 recomputed with mbedTLS 40 of 40
authentication tags, under two keys per session separated by a
zeroization \meas{}; the rejection of a divider below the minimum and of
one above the maximum was also exercised on the board. Known-answer
vectors cannot be loaded through the register map, which accepts neither
a key nor a seed, and remain covered in simulation. Getting there took
corrections that no simulation had asked for --- a ring that would not
route, a ring domain and a health-test comparison that did not close
timing, and inverters that the tool had merged --- and the rule
learned on the first prototype board: check in the implemented netlist
that the ring exists.

\paragraph{The rings of the FPGA have a measurable white noise floor, and
little else in their bits is entropy (O6).} The estimators of the
relative phase noise give between $3\times10^{-6}$ and $8\times10^{-5}$
per period of the sampling ring \meas{} depending on the method and on
the logic attached to the rings, each repeatable within a few per cent:
the stochastic model does not describe the shape of what was measured,
and the accumulation exponent is 1.6 to 2.1. A spectrum explains most of
it: a narrow line near 39.4~kHz, of external and unidentified origin,
modulates the relative phase. Away from the line the white floor is 0.6
to $1.1\times10^{-6}$ \meas{}, with an estimator that errs upwards, and
its lower edge puts the \ac{FPGA} divider at \kd{19}, for
$H_\infty \ge 0.99$ in the model at about 390~bit/s. The divider first
derived from the fits, $2^{16}$, would have given 0.47. Two findings bear
on the design rather than on the board. The health tests of SP~800-90B do
not fire when the divider is lowered until the bits simply alternate, so
a minimum divider enforced in hardware is the only protection of that
regime. And the logic attached to a ring changes its apparent phase noise
several times over, which is why the integrated circuit asks for
isolated rings. Three steps of the protocol fixed beforehand were not
followed as written, and Section~\ref{sec:campaign} declares each.

\paragraph{Comparison of the two sources (O8).} Under the SP~800-90B
battery the assigned figure does not separate the \ac{FPGA} source
(0.857), the ESP32 (0.848) and ideal bits (0.847) at $10^6$ bits
\meas{}. The estimators sensitive to bias do separate the \ac{FPGA}, by a
bias of half a per cent that follows the implementation, and so the
battery gives the better marks to the generator that, without boot
noise, repeats its sequence. At low dividers the battery overestimates,
as the simulation predicted, instead of contradicting the model. What
separates the two sources is the restart test, which the ring source
passes with nothing but its own noise in 1000 reconfigurations and the
ESP32 fails without the noise injected at boot; and, behind the figures,
a physical quantity measured inside the device, which the ESP32 does not
offer. The long records of the \ac{FPGA} were taken at a divider that
the white floor later showed to be too low, and they say something about
the battery, not about the source.

\section{Contributions}

\begin{itemize}
\item A chain from transistor noise to a min-entropy bound for a ring
oscillator on an open \ac{PDK}, with the divider fixed by the worst of
corners, supply and mismatch, and with every link validated on data of
known answer.
\item A Monte Carlo of the sampling flip-flop that turned its bias into a
design decision, the length of the sampled ring, under a criterion fixed
before the simulation.
\item Evidence, on simulated sources of known entropy and then on a real
one, that the SP~800-90B estimators saturate below 0.9 at the usual
sequence lengths and cannot tell a sound source from a weaker one in that
range.
\item The white floor of the relative phase noise of two rings in an
Artix-7, measured inside the device and separated from a deterministic
modulation that the usual fits count as jitter, with the limits of the
estimator stated; and the observation that the standard health tests are
blind to the anticorrelated low-jitter regime, which is the case for a
minimum divider enforced in hardware.
\item A restart test of the \ac{FPGA} source set beside that of the ESP32
generator.
\item A characterization of the ESP32 generator under SP~800-90B with
three chips and four conditions, and a restart test showing that, without
external noise at boot, different chips produce the same sequence. This
result was not found in the literature consulted.
\item A key-generation engine in portable RTL with its security rules
written as checked properties, synthesized on open standard cells, with
its firmware and reproducible flow released.
\end{itemize}

\section{Future work}

\begin{itemize}
\item \emph{Silicon.} Fabricate the ring and the sampler on an academic
shuttle and measure the jitter that this work simulated; the \ac{FPGA}
measures a different ring.
\item \emph{Foundry models.} Repeat the corner and mismatch study with the
noise models of a commercial \ac{PDK}.
\item \emph{What the simulation of the final rings leaves open.} Mismatch
with eleven stages and the jitter of the thirteen-stage ring; the timing
of the prescaled domain with a slow library; and a sampling clock drawn
at random in the flip-flop Monte Carlo, to settle the difference between
its batches.
\item \emph{The line near 39.4~kHz.} Identify its origin, with an
oscilloscope on the supplies of the board or with the ESP32 isolated, and
a model of the shape of the measured autocorrelation that includes it.
\item \emph{What the campaign did not do.} A long record at the operating
divider, the rings in adjacent regions, and temperature and supply on the
board. The stream is limited by the serial link of the ESP32; storing the
string before sending it would extend it to lower dividers.
\item \emph{The meter with the real engine.} The signed-readings firmware
ran against a model of the engine; run it against the \ac{FPGA}.
\item \emph{Online jitter test.} Move the pair estimator into the RTL as an
online test that detects cooling and locking, as its authors propose.
\item \emph{Area.} Apply the third optimization to approach 45~kGE, and
measure power.
\item \emph{The product.} Store the device key encrypted or make the
engine a secure element that signs, and add real time to the signed
readings.
\item \emph{Formal verification} of the sequencer properties, once the
tool chain reads them.
\end{itemize}
